\chapter{The Multiscale Graph VFE and Closing the Tower}
\label{ch:full-graph-meta-agent-vfe}

This chapter asks what probabilistic structure places adaptive event laws and links
on the interaction multigraph, random partitions, and the influence of parents
on their children inside one variational free energy, and how a tower of such depths
closes. The answer is the free energy of one normalized generative law and one
recognition law over the whole tower (\Cref{def:full-graph-tower-vfe}). For an
adaptive tower both laws factor top-down, each depth conditioned on the whole
next-depth state and the structural mark (\Cref{lem:full-graph-adaptive-reversal}),
and the free energy splits into observation, top-prior, partition, graph, and
parent--child terms (\Cref{thm:full-graph-vfe-decomposition}). A finite acyclic tower
needs a declared latent root (\Cref{prop:obs-declared-root-unavoidable});
participation across depths enters through the inference loop, parameter learning
(\Cref{prop:full-graph-aggregate-prior}), and delayed feedback
(\Cref{prop:full-graph-episode-joint,prop:full-graph-episode-lifted-joint}). The formation witnesses of
\Cref{ch:finite-depth-formation} are adaptive towers, and \Cref{ch:philosophy} reads
the three loci of participation.

The depth $\ell$ runs from $0$ to a finite $L$. At depth $\ell$, $Z_\ell$ are the
coordinates, $\mathfrak N_\ell$ the graph datum, $Y_\ell=(Z_\ell,\mathfrak N_\ell)$ the
state on a standard Borel space $\mathsf Y_\ell$ that here includes the graph datum,
$R_\ell\in\mathsf R_\ell$ the structural mark with realized value $\gamma$, and
$\mathfrak h_\ell$ the retained root, holonomy, and boundary marks. The whole tower is
$\mathbf Y$, row normalizers are $\mathcal Z$, and $\mathbb M$ is an input law:
generative, posterior, or recognition. In the episode-indexed models of
\Cref{sec:obs-participatory}, $\widetilde o_n$ is a latent copy of the record $o_n$,
$\epsilon_{\mathrm e}$ is the constant step of the aggregate parameter update, and
$\mathrm T$ is the temperature of the descent kinetics, the symbol of
\Cref{ch:finite-depth-formation}.

Scope. In an adaptive tower the downward kernel is a reversal kernel, and the coupling
it carries is not read as causal downward influence; declining that reading does not
refute a causal reading of a separately declared top-down hierarchy. The composite
potential of \Cref{def:full-graph-composite-potential} is not claimed to be an
evidence bound, and delayed feedback is not claimed to support a sustained
nonequilibrium reading. \status{NOT-CLAIMED} Open items are the persistent hierarchy
(\Cref{open:full-graph-hierarchy}), the accounting across episodes
(\Cref{open:full-graph-episode-accounting}), the dynamic closure of the pointwise
parent (\Cref{open:rg-pointwise-parent-dynamics}), and a natural partition selector
(\Cref{open:cg-natural-selector}).

\section{The full answer before decomposition}
\label{sec:full-graph-vfe-global}

\definitionheading{Tower variational free energy}{def:full-graph-tower-vfe}
Write $X_*=X(c_*)$, fix an admitted observation $o\in\mathsf
O^{\mathrm{reg}}_{\theta,X_*}$ and a finite depth $L$, and let $\mathbf Y$ contain
every random object of the multiscale model: belief and model coordinates at depths
$0,\ldots,L$, links, event laws, retained marks, structural marks, and auxiliary
latents. Declare one normalized generative kernel $\mathbb
P_\theta(do,d\mathbf Y\given X_*)$ and one recognition kernel $\mathbb
Q_\lambda(d\mathbf Y\given o,X_*)$. With posterior slice
$\boldsymbol\Pi_{\theta,o,X_*}$ and evidence density $p_\theta(o\given
X_*)\in(0,\infty)$, the tower variational free energy is
\begin{equation}
\boxed{
\Fenergy^{\mathrm{ext}}_{\mathrm{tower}}[\mathbb Q_\lambda;X_*,o]
=-\log p_\theta(o\given X_*)
+\KL\bigl(\mathbb Q_\lambda(\cdot\given o,X_*)\Vert\boldsymbol\Pi_{\theta,o,X_*}\bigr).}
\label{eq:full-graph-vfe-global}
\end{equation}
By the evidence identity of \Cref{ch:elbo}, on the finite-density domain it equals
\begin{equation}
\Fenergy_{\mathrm{tower}}
=\E_{\mathbb Q_\lambda}\Bigl[\log\frac{q_\lambda(\mathbf Y\given o,X_*)}{p_\theta(o,\mathbf Y\given X_*)}\Bigr].
\label{eq:full-graph-vfe-density}
\end{equation}
Every later term is an accounting decomposition of this object or is named a
composite potential.

\section{A normalized hierarchical graph model}
\label{sec:full-graph-vfe-joint}

\definitionheading{Closed scale datum}{def:full-graph-scale-datum}
At depth $\ell$ let $V_\ell$ be the vertex set and let $Z_\ell=(Z_{i,\ell})_{i\in
V_\ell}$ contain the full probabilistic coordinates of that depth, belief and model
presentations rather than only their marginals. The graph datum lives on a finite
interaction multigraph $\mathfrak I_\ell=(V_\ell,E_\ell,\bar{\ })$
(\Cref{def:geo-graph-links}), at depth $0$ the multigraph of \cref{sec:intro-conventions} and at deeper depths
induced by the link lift of \Cref{def:full-graph-coarse-links}. Put
\begin{equation}
\eta_\ell=(\eta^{b,\ell},\eta^{m,\ell}),
\qquad
\Theta^{x,\ell}=(\Theta^{x,\ell}_e)_{e\in E_\ell},
\qquad
\mathfrak N_\ell=(\eta_\ell,\Theta^{b,\ell},\Theta^{m,\ell},\mathfrak h_\ell),
\label{eq:full-graph-closed-scale-datum}
\end{equation}
where $\eta^{x,\ell}=(\eta^{x,\ell}_e)_{e\in E_\ell}$ is a normalized event law on the
copies, $\Theta^{x,\ell}_e\in G$ is the channel-$x$ link of the copy $e$ with
$\Theta^{x,\ell}_{\bar e}=(\Theta^{x,\ell}_e)^{-1}$, and $\mathfrak h_\ell$ contains
the root, based-holonomy, and dressed boundary marks of
\Cref{sec:rg-gauge-cross-scale}. Each copy carries one link, one event mass, and one
divergence per channel; endpoint notation abbreviates a copy only when it is unique,
and $\eta^{x,\ell}_{ij}$ is the total mass of the copies $j\to i$. Occupancies and
rows are derived from $\eta_\ell$ as in \eqref{eq:rg-attention-event-law} and do not
replace it in the recursive state. The observation and $X_*$ stay outside every
coarse channel.

\definitionheading{Common structural channel}{def:full-graph-structural-channel}
A structural step is one recognition-independent Markov channel
\begin{equation}
C^{\mathrm{struct}}_{\ell+1\leftarrow\ell}:
\mathsf Y_\ell\rightsquigarrow\mathsf R_\ell\times\mathsf Y_{\ell+1},
\label{eq:full-graph-common-structural-channel}
\end{equation}
whose output includes $(R_\ell,\eta_{\ell+1},\mathfrak h_{\ell+1})$, with augmented
transition state
\begin{equation}
\widehat{\mathsf Y}_{\ell+1\leftarrow\ell}
:=(\mathsf Y_\ell,R_\ell,\mathsf Y_{\ell+1}),
\label{eq:full-graph-augmented-transition-state}
\end{equation}
which does not duplicate $R_\ell$. The same channel is applied to the fine generative
marginal, the selected posterior, and the recognition law,
\begin{equation}
\widehat{\mathbb M}^{\mathrm{struct}}_{\ell+1\leftarrow\ell}(dy_\ell,d\gamma,dy_{\ell+1})
:=\mathbb M_\ell(dy_\ell)C^{\mathrm{struct}}_{\ell+1\leftarrow\ell}(y_\ell;d\gamma,dy_{\ell+1}),
\qquad
\mathbb M\in\{\mathbb P,\boldsymbol\Pi,\mathbb Q\},
\label{eq:full-graph-common-structural-pushforwards}
\end{equation}
so the three laws of $R_\ell$ may differ while the channel does not. If the channel is
realized by an auxiliary seed, the same measurable map and seed law define all three
pushforwards; a blocking rule that depends on the input law is not admitted.

\hypothesisheading{Depth-Markov hierarchy}{hyp:full-graph-depth-markov}
The generative law factors top-down as
\begin{align}
\mathbb P_\theta(do,d\mathbf Y\given X_*)
={}&\mathbb P^{\mathrm{obs}}_\theta(do\given Z_0,X_*)\mathbb P_L(dY_L\given X_*)
\nonumber\\
&\times\prod_{\ell=0}^{L-1}
\mathbb P^\ell_R(dR_\ell\given Y_{\ell+1},X_*)
\mathbb P^\ell_{\mathfrak N}(d\mathfrak N_\ell\given Y_{\ell+1},R_\ell,X_*)
K^\ell_\downarrow(dZ_\ell\given Y_{\ell+1},R_\ell,\mathfrak N_\ell,X_*),
\label{eq:full-graph-generative-factorization}
\end{align}
with the kernels composed from depth $L$ downward, and recognition factors along the
same order,
\begin{align}
\mathbb Q_\lambda(d\mathbf Y\given o,X_*)
={}&\mathbb Q_L(dY_L\given o,X_*)\prod_{\ell=0}^{L-1}
\mathbb Q^\ell_R(dR_\ell\given Y_{\ell+1},o,X_*)
\nonumber\\
&\times\mathbb Q^\ell_{\mathfrak N}(d\mathfrak N_\ell\given Y_{\ell+1},R_\ell,o,X_*)
\mathbb Q^\ell_Z(dZ_\ell\given Y_{\ell+1},R_\ell,\mathfrak N_\ell,o,X_*).
\label{eq:full-graph-recognition-factorization}
\end{align}
Each depth-$\ell$ factor reads the complete next-depth state
$Y_{\ell+1}=(Z_{\ell+1},\mathfrak N_{\ell+1})$ and no deeper variable; this excludes
direct dependence on depths $\ell+2,\ldots,L$ given $(Y_{\ell+1},R_\ell)$, leaves
arbitrary dependence within a depth and between adjacent depths, and is not a
mean-field assumption. In a declared top-down hierarchy
$K^\ell_\downarrow$ is where downward influence enters, and the bottom-up coarse map
belongs to recognition or to an induced pushforward, never to a second generative
arrow at the same depth.

\lemmaheading{Depth-Markov reversal of an adaptive tower}{lem:full-graph-adaptive-reversal}
Let $\mathbb M_0$ be a probability law on $\mathsf Y_0$, and generate the adaptive
tower bottom-up through the common structural channels,
\begin{equation}
\mathbb M_0(dy_0)\prod_{\ell=0}^{L-1}
C^{\mathrm{struct}}_{\ell+1\leftarrow\ell}(y_\ell;d\gamma_\ell,dy_{\ell+1}).
\label{eq:full-graph-adaptive-tower-law}
\end{equation}
For every $\ell$, the conditional law of $R_\ell$ given all variables
$(Y_{\ell+1},R_{\ell+1},\ldots,Y_L)$ of greater depth has a version that depends only
on $Y_{\ell+1}$, and the conditional law of $Y_\ell$ given $R_\ell$ and those variables
has a version that depends only on $(R_\ell,Y_{\ell+1})$. Hence
\eqref{eq:full-graph-adaptive-tower-law} factors as in
\eqref{eq:full-graph-generative-factorization} without its observation factor, with
$\mathbb P^\ell_R$, $\mathbb P^\ell_{\mathfrak N}$, and $K^\ell_\downarrow$ selected
versions of these conditional laws. The same holds for the recognition input $\mathbb
Q_0(dy_0\given o,X_*)$, for the selected posterior slice, and for the generative input
$\mathbb P_0(do,dy_0\given X_*)$ whenever its observation kernel reads only $Z_0$, so
\Cref{hyp:full-graph-depth-markov} holds automatically for an adaptive tower. These
factors are reversal, or bridge, kernels in the sense of
\eqref{eq:rg-agent-meta-bayes-adjoint}, not independently chosen structural
channels, and when the parent coordinates are a function of the child coordinates and
the structural mark, $K^\ell_\downarrow$ is supported on the child configurations
compatible with the realized parent. \status{ESTABLISHED}

\noindent\emph{Proof.}
For each $\ell$, the variables $(R_{\ell+1},Y_{\ell+2},\ldots,Y_L)$ are produced from
$Y_{\ell+1}$ by kernels that read neither the observation nor any of
$(Y_0,R_0,\ldots,Y_\ell,R_\ell)$, so they are conditionally independent of the latter
given $Y_{\ell+1}$, which gives both version statements. Regular conditional laws
exist on standard Borel spaces \citep{Kallenberg2021}, and successive disintegration
from depth $L$ downward gives the factorization. A conditional law of $Z_\ell$ given a
value of a function of $(Z_\ell,R_\ell)$ charges only the preimage of that value.
$\square$

\propositionheading{The next-depth graph datum cannot be dropped}{prop:full-graph-dropped-datum}
If the depth-$\ell$ factors of an adaptive tower read only $(Z_{\ell+1},R_\ell)$, the
product $\mathbb P^{\mathrm{red}}$ of the corresponding selected conditional laws of
the tower law $\mathbb P^{\mathrm{tow}}$ of \eqref{eq:full-graph-adaptive-tower-law}
differs from it by
\begin{equation}
\KL\bigl(\mathbb P^{\mathrm{tow}}\Vert\mathbb P^{\mathrm{red}}\bigr)
=\sum_{\ell=0}^{L-1}
I_{\mathbb P^{\mathrm{tow}}}\bigl(Y_\ell;\mathfrak N_{\ell+1}\given Z_{\ell+1},R_\ell\bigr).
\label{eq:full-graph-dropped-datum-cost}
\end{equation}
The cost vanishes exactly when every $\mathfrak N_{\ell+1}$ is conditionally
independent of $Y_\ell$ given $(Z_{\ell+1},R_\ell)$. For a deterministic event
pushforward it is positive as soon as two event laws in the support, at the same
$(Z_{\ell+1},R_\ell)$, are pushed to different coarse event laws: for four fine
vertices with the partition $\{0,1\}\mapsto I$, $\{2,3\}\mapsto J$, constant
coordinates, trivial links and marks, and $\eta_0$ equal with probability $1/2$ each
to the unit mass on a copy $2\to0$ or on a copy $1\to0$, the pushed $\eta_1$ is the
unit mass on $(I,J)$ or on $(I,I)$, and the divergence is $\log2$.
\status{ESTABLISHED}

\noindent\emph{Proof.}
By \Cref{lem:full-graph-adaptive-reversal} the reduced product replaces each
conditional law of $Y_\ell$ given $(R_\ell,Y_{\ell+1})$ by its conditional law given
$(R_\ell,Z_{\ell+1})$, and the chain rule \Cref{lem:prob-kl-chain-rule} turns each
replacement into the displayed conditional mutual information. In the witness the
joint law of $(\mathfrak N_0,\mathfrak N_1)$ has two atoms of mass $1/2$ while the
reduced product is the product of their marginals, which puts mass $1/2$ on pairs
that violate the pushforward. $\square$

\theoremheading{Exact hierarchical decomposition}{thm:full-graph-vfe-decomposition}
Assume \Cref{hyp:full-graph-depth-markov}, which holds for an adaptive tower by
\Cref{lem:full-graph-adaptive-reversal}, together with common domination and the
integrability needed for the finite log-density form. Then
\begin{align}
\Fenergy_{\mathrm{tower}}
={}&-\E_{\mathbb Q_\lambda}\log p^{\mathrm{obs}}_\theta(o\given Z_0,X_*)
+\KL(\mathbb Q_L\Vert\mathbb P_L)
\nonumber\\
&+\sum_{\ell=0}^{L-1}\E_{\mathbb Q_\lambda}\bigl[
\KL(\mathbb Q^\ell_R\Vert\mathbb P^\ell_R)
+\KL(\mathbb Q^\ell_{\mathfrak N}\Vert\mathbb P^\ell_{\mathfrak N})
+\KL(\mathbb Q^\ell_Z\Vert K^\ell_\downarrow)\bigr],
\label{eq:full-graph-vfe-factorized}
\end{align}
where each divergence compares the recognition and generative factors of
\eqref{eq:full-graph-recognition-factorization} and
\eqref{eq:full-graph-generative-factorization} at their common conditioning
variables, the recognition factor also reading $o$, and is averaged over those
variables under $\mathbb Q_\lambda$. These are the observation, top-prior,
partition, graph, and parent--child contributions. \status{ESTABLISHED}

\noindent\emph{Proof.}
Substitute the two factorizations into the Radon--Nikodym ratio of
\eqref{eq:full-graph-vfe-density} and apply the chain rule for conditional relative
entropy, \Cref{lem:prob-kl-chain-rule} and \citet[Sec.~2.5]{CoverThomas2006}, in its
measure-level form, which subtracts no infinities. $\square$

\section{The scale-local two-channel graph energy}
\label{sec:full-graph-vfe-scale-local}

\definitionheading{Transported divergences and row free energies}{def:full-graph-transported-divergences}
For every copy $e:j\to i$ in $E_\ell$ and $x\in\{b,m\}$ put
\begin{equation}
D^{x,\ell}_e:=\KL\bigl(q^x_{i,\ell}\Vert(\Theta^{x,\ell}_e)_\#q^x_{j,\ell}\bigr),
\label{eq:full-graph-transported-divergences}
\end{equation}
where $(\Theta)_\#$ is the pushforward action $\widehat\rho_x(\Theta)$ on channel-$x$
laws; parallel copies carry separate divergences, and these are law-point energies in
the sense of \Cref{sec:prob-levels}. A selectively communicating variant uses the
role divergences \eqref{eq:obs-selective-role-divergence} or a message channel
instead, which changes the energy, and its recursive state keeps the interface and
noise data and any private conditional law or memory that later kernels read
(\Cref{sec:selective-records}, \Cref{prop:geo-selective-descent}). With $E_\ell(i)$
the copies received by $i$, the occupancies and rows of an event law are
\begin{equation}
\alpha^{x,\ell}_i:=\sum_{e\in E_\ell(i)}\eta^{x,\ell}_e,
\qquad
\beta^{x,\ell}_e:=\frac{\eta^{x,\ell}_e}{\alpha^{x,\ell}_i}
\quad(e\in E_\ell(i),\ \alpha^{x,\ell}_i>0),
\label{eq:full-graph-edge-event-disintegration}
\end{equation}
with pair rows $\beta^{x,\ell}_{ij}=\sum_{e:j\to i}\beta^{x,\ell}_e$. The rows are the
conditional source laws of whichever event law is displayed, generative, posterior,
or recognition; a row of zero occupancy is not determined by the joint law. With
normalized copy priors $\pi^{x,\ell}_i$ on $E_\ell(i)$ at every vertex $i$ that carries a row (a
vertex with an empty source set carries no row and has zero occupancy, although $E_\ell(i)$ may
contain the reversals of copies it sends) and temperatures
$\tau^x_{i,\ell}>0$, the row free energies in energy units are
\begin{equation}
\Fenergy^{\mathrm{att},x}_{i,\ell}
=\sum_{e\in E_\ell(i)}\beta^{x,\ell}_eD^{x,\ell}_e
+\tau^x_{i,\ell}\KL\bigl(\beta^{x,\ell}_i\Vert\pi^{x,\ell}_i\bigr),
\label{eq:full-graph-row-free-energies}
\end{equation}
and their dimensionless form is $\Fenergy^{\mathrm{att},x}_{i,\ell}/\tau^x_{i,\ell}$.

\propositionheading{Gibbs rows}{prop:full-graph-gibbs-rows}
With $\mathcal Z^{x,\ell}_i=\sum_{e\in E_\ell(i)}\pi^{x,\ell}_e
e^{-D^{x,\ell}_e/\tau^x_{i,\ell}}$, assume $\mathcal Z^{x,\ell}_i>0$, equivalently that at least
one edge of positive prior mass has finite energy. On the closed simplex of rows supported on
$\{e:\pi^{x,\ell}_e>0\}$, with $0\cdot(+\infty)=0$, the extended-real identity is
$\Fenergy^{\mathrm{att},x}_{i,\ell}=\tau^x_{i,\ell}\KL(\beta^{x,\ell}_i\Vert\beta^{x,\ell,*}_i)
-\tau^x_{i,\ell}\log\mathcal Z^{x,\ell}_i$ with
\begin{equation}
\beta^{x,\ell,*}_e=\frac{\pi^{x,\ell}_ee^{-D^{x,\ell}_e/\tau^x_{i,\ell}}}{\mathcal Z^{x,\ell}_i},
\label{eq:full-graph-gibbs-rows}
\end{equation}
so the minimizer is unique and the minimum is $-\tau^x_{i,\ell}\log\mathcal
Z^{x,\ell}_i$, with zero optimal mass on every infinite-energy edge. If $\mathcal Z^{x,\ell}_i=0$,
every prior-supported row has infinite free energy, and the Gibbs formula defines no probability row
or finite minimizer. When all prior-supported energies are finite, the optimum has positive mass on
every prior-supported edge. The optimal rows move with the transported divergences, and omitting
the row relative entropy changes the objective. The dimensionless form is the row
free energy \eqref{eq:obs-attention-vfe} with prior $\pi^{x,\ell}_i$ and the fixed
energies $D^{x,\ell}_e$ in place of the expected divergences, with the same minimizer
\eqref{eq:obs-attention-recognition-optimum} and minimum $-\log\mathcal
Z^{x,\ell}_i$. The energy-unit form with $\tau^x_{i,\ell}\ne1$ is not an
independently weighted sector of one evidence bound (\Cref{prop:obs-attention-elbo}),
and the composite potential below inherits this scaling. \status{ESTABLISHED}

\noindent\emph{Proof.}
Expanding $\KL(\beta\Vert\beta^*)$ with \eqref{eq:full-graph-gibbs-rows} gives
$\sum_e\beta_e\log(\beta_e/\pi_e)+\sum_e\beta_eD_e/\tau+\log\mathcal Z$, which is the
stated identity after multiplication by $\tau$. Relative entropy is nonnegative and
vanishes only at $\beta=\beta^*$. A row charging an infinite-energy edge has both sides equal to
$+\infty$; otherwise the expansion uses only finite-energy edges. If $\mathcal Z=0$, every edge of
positive prior mass has infinite energy, giving the stated exceptional case. $\square$

\definitionheading{Composite multiscale potential}{def:full-graph-composite-potential}
If a scale-local prior is a normalized Gibbs law containing the self terms and
$\sum_{e\in E_\ell}(\eta^{b,\ell}_eD^{b,\ell}_e+\eta^{m,\ell}_eD^{m,\ell}_e)$, those
terms and the global normalizer belong to the tower variational free energy. If they
are appended to a fixed population free energy $\Fenergy_{\mathrm{pop}}$ of
\Cref{ch:local-collective-elbo} instead, the resulting objective is the composite
potential
\begin{equation}
\mathcal J_{\mathrm{comp}}
=\Fenergy_{\mathrm{pop}}
+\sum_\ell\sum_i\bigl(\alpha^{b,\ell}_i\Fenergy^{\mathrm{att},b}_{i,\ell}
+\alpha^{m,\ell}_i\Fenergy^{\mathrm{att},m}_{i,\ell}\bigr)
+\sum_\ell\Fenergy^\ell_{\mathrm{part}}
+\sum_\ell\Fenergy^\ell_{\mathrm{cross}}
+\sum_\ell\mathcal J^\ell_{\mathrm{link}},
\label{eq:full-graph-composite-potential}
\end{equation}
where the sum over $i$ runs over the vertices that carry a row, a term of zero occupancy is
zero, the partition and cross-scale terms are those of \eqref{eq:full-graph-partition-free-energy}
and \eqref{eq:full-graph-cross-scale-kl}, and $\mathcal J^\ell_{\mathrm{link}}$ is a declared link
regularizer. It is a descent potential and carries no evidence
interpretation of its own.

\section{Partition variables and event pushforwards}
\label{sec:full-graph-vfe-partitions}

Attention rows do not create a parent random variable. A hierarchy meant to emerge
under the same variational principle puts the structural mark $R_\ell$ inside both
the generative and the recognition law.

\definitionheading{Partition factor}{def:full-graph-partition-factor}
A finite-temperature partition term has the form
\begin{equation}
\Fenergy^\ell_{\mathrm{part}}
=\E_{\mathbb Q^\ell_R}\bigl[H^\ell_{\mathrm{block}}(R_\ell;Z_\ell,\mathfrak N_\ell)\bigr]
+\tau^\ell_R\KL\bigl(\mathbb Q^\ell_R\Vert\mathbb P^{\ell,\circ}_R\bigr),
\label{eq:full-graph-partition-free-energy}
\end{equation}
and enters the variational free energy only when the normalized Gibbs factor that it
represents is part of the generative law. The block action
$H^\ell_{\mathrm{block}}$ may trade internal transported mismatch against boundary
retention, evaluator closure, and residual cost. No canonical partition selector
follows from the global free energy alone (\Cref{open:cg-natural-selector}).

\subsection{A random partition catalog}
\label{sec:rg-geometry-normalization}

The structural mark may be random with a finite catalog of partitions, or membership
configurations, of $V_\ell$; each label fixes the branch carrier of the next depth,
$X$ stays outside the channel, no edge is created, and the two normalizations of the
catalog are distinct generative models.

\definitionheading{Finite partition catalog}{def:rg-geometry-ensemble}
Let $\mathfrak G$ be a finite nonempty catalog of values of $R_\ell$. For each
$\gamma\in\mathfrak G$ declare a finite nonempty carrier $\mathsf Y_\gamma$, a
strictly positive probability reference $\nu_\gamma$, a positive raw factor
$\widetilde p_\gamma$, a positive weight $\mathsf a_\gamma$, and a normalized,
possibly vanishing, likelihood $L_\gamma(o\given y)$ on a common finite observation
alphabet, all fixed generative data independent of recognition. On the disjoint union
of the carriers put
\begin{equation}
\mathcal Z_\gamma=\sum_y\nu_\gamma(y)\widetilde p_\gamma(y),
\qquad
p_\gamma(y)=\frac{\nu_\gamma(y)\widetilde p_\gamma(y)}{\mathcal Z_\gamma},
\qquad
\bar{\mathsf a}_\gamma=\frac{\mathsf a_\gamma}{\sum_\delta\mathsf a_\delta}.
\label{eq:rg-geometry-branch}
\end{equation}
The hierarchically normalized and the globally normalized joints are
\begin{align}
\mathbb P_{\mathrm h}(\gamma,y,o)
&=\bar{\mathsf a}_\gamma p_\gamma(y)L_\gamma(o\given y),
\label{eq:rg-geometry-hierarchy}\\
\mathbb P_{\mathrm f}(\gamma,y,o)
&=\frac{\mathsf a_\gamma\nu_\gamma(y)\widetilde p_\gamma(y)L_\gamma(o\given y)}{\mathcal Z_{\mathrm{all}}},
\qquad
\mathcal Z_{\mathrm{all}}=\sum_\delta\mathsf a_\delta\mathcal Z_\delta .
\label{eq:rg-geometry-factor}
\end{align}

\propositionheading{Branch normalization changes label weights}{prop:rg-geometry-weights}
Both joints are normalized and have the same conditional branch laws, while the label
weight of the globally normalized joint is $\mathsf a_\gamma\mathcal
Z_\gamma/\mathcal Z_{\mathrm{all}}$. The label weights of the two joints are the
softmax weights of \Cref{def:rg-thermo-finite-score} with effective log masses
$\log\mathsf a_\gamma$ and $\log\mathsf a_\gamma+\log\mathcal Z_\gamma$
respectively. For the same positive weights the joints coincide exactly when all
$\mathcal Z_\gamma$ coincide. Replacing $\widetilde p_\gamma$ by $\mathrm
c_\gamma\widetilde p_\gamma$ with $\mathrm c_\gamma>0$ leaves the hierarchical joint
unchanged and reweights the global one by $\mathrm c_\gamma$, and a common factor
cancels in both. For a strictly positive probability reference $\nu_{\mathfrak G}$ on
the catalog write $\mathsf a_\gamma=\nu_{\mathfrak G}(\gamma)e^{-H^{\mathrm{lab}}_\gamma}$ and
$\widetilde p_\gamma=e^{-H_\gamma}$; every positive family of weights has this form.
Fix positive evidence $M_o=\sum_{\gamma,y}\mathbb P(\gamma,y,o)$ and a recognition law
$\mathbb Q(\gamma,y)=\mathbb Q^{\mathrm{lab}}(\gamma)\mathbb Q_\gamma(y)$, with
$\mathbb Q_\gamma$ defined where $\mathbb Q^{\mathrm{lab}}(\gamma)>0$. Then, with
$\Fenergy_o=\KL(\mathbb Q\Vert\boldsymbol\Pi_o)-\log M_o$,
\begin{align}
\Fenergy_{\mathrm f,o}(\mathbb Q)
&=\KL(\mathbb Q^{\mathrm{lab}}\Vert\nu_{\mathfrak G})+\E_{\mathbb Q^{\mathrm{lab}}}H^{\mathrm{lab}}_\gamma
+\log\mathcal Z_{\mathrm{all}}
\nonumber\\
&\quad+\sum_\gamma\mathbb Q^{\mathrm{lab}}(\gamma)
\bigl[\KL(\mathbb Q_\gamma\Vert\nu_\gamma)+\E_{\mathbb Q_\gamma}H_\gamma
-\E_{\mathbb Q_\gamma}\log L_\gamma(o\given Y)\bigr],
\label{eq:rg-geometry-vfe-f}\\
\Fenergy_{\mathrm h,o}(\mathbb Q)
&=\Fenergy_{\mathrm f,o}(\mathbb Q)-\log\mathcal Z_{\mathrm{all}}
+\log\sum_\delta\mathsf a_\delta
+\sum_\gamma\mathbb Q^{\mathrm{lab}}(\gamma)\log\mathcal Z_\gamma,
\label{eq:rg-geometry-vfe-h}
\end{align}
with the convention $0\log0=0$; a recognition-positive atom with zero likelihood
gives $+\infty$. Each expansion contains one full recognition entropy, decomposed by
its conditional chain, and no further correlation-entropy term.
\status{ESTABLISHED}

\noindent\emph{Proof.}
Sum either joint over $o$, then $y$, then $\gamma$. Equality of the joints implies
equality of the label marginals, hence $\mathcal Z_\gamma=\mathcal
Z_{\mathrm{all}}/\sum_\delta\mathsf a_\delta$ for every $\gamma$; the converse is
substitution. A branch rescaling changes $\mathcal Z_\gamma$ to $\mathrm
c_\gamma\mathcal Z_\gamma$ and cancels in $p_\gamma$. The softmax statement is the
label marginal of each joint written as the exponential of its log mass divided by the
sum of these exponentials. The expansions follow by writing $\log\mathbb P$ with the
displayed factors and taking the expectation under $\mathbb Q$; the two differ only in
the normalizers, $\log\sum_\delta\mathsf a_\delta+\log\mathcal Z_\gamma$ against
$\log\mathcal Z_{\mathrm{all}}$. $\square$

Reference, raw factor, label weight, evidence, and posterior are distinct objects,
and the score of either joint in a parameter is the complete labeled score of
\Cref{prop:rg-thermo-finite-labeled-fisher}, which contains every varying generative
factor, the $\mathcal Z_\gamma$ of the global model included.

\subsection{Membership kernels and event pushforwards}

\definitionheading{Membership kernels}{def:full-graph-membership-kernel}
Overlapping assignment is a membership Markov kernel
\begin{equation}
\mathfrak m_\ell(I\given i,R_\ell)\ge0,
\qquad
\sum_{I\in V_{\ell+1}}\mathfrak m_\ell(I\given i,R_\ell)=1,
\label{eq:full-graph-membership-kernel}
\end{equation}
of which a hard partition is the deterministic case $\mathfrak m_\ell(I\given
i,R_\ell)=\mathbf 1\{R_\ell(i)=I\}$. Correlated endpoints use one joint endpoint
kernel
\begin{equation}
\mathfrak M_\ell(I,J\given i,j,R_\ell)\ge0,
\qquad
\sum_{I,J}\mathfrak M_\ell(I,J\given i,j,R_\ell)=1,
\label{eq:full-graph-joint-endpoint-kernel}
\end{equation}
whose product form requires a declared conditional independence of the two draws.

\propositionheading{Event pushforwards}{prop:full-graph-event-pushforward}
Fix $y\in\mathsf Y_\ell$ and a realized mark $\gamma\in\mathsf R_\ell$. If every
vertex receives one membership draw and the draws of distinct vertices are
conditionally independent at $(y,\gamma)$, the conditional coarse event law is
\begin{equation}
\eta^{x,\ell+1}_{IJ}(y,\gamma)
=\sum_{\substack{i,j\in V_\ell\\ i\ne j}}\eta^{x,\ell}_{ij}(y)
\mathfrak m_\ell(I\given i,\gamma)\mathfrak m_\ell(J\given j,\gamma)
+\delta_{IJ}\sum_{i\in V_\ell}\eta^{x,\ell}_{ii}(y)\mathfrak m_\ell(I\given i,\gamma),
\label{eq:full-graph-edge-pushforward}
\end{equation}
where the loop term uses the diagonal endpoint kernel, since both endpoints of a loop
follow one draw. With correlated endpoints,
\begin{equation}
\eta^{x,\ell+1}_{IJ}(y,\gamma)
=\sum_{i,j\in V_\ell}\eta^{x,\ell}_{ij}(y)\mathfrak M_\ell(I,J\given i,j,\gamma).
\label{eq:full-graph-correlated-edge-pushforward}
\end{equation}
With the assignment marginal $C^R_\ell(y,d\gamma)=C^{\mathrm{struct}}_{\ell+1\leftarrow\ell}(y;d\gamma\times\mathsf
Y_{\ell+1})$ of the structural channel and an input law $\mathbb M$, the unconditional
event mass averages the whole conditional expression,
\begin{equation}
\overline\eta^{x,\ell+1}_{IJ}[\mathbb M]
=\sum_{i,j\in V_\ell}\int\eta^{x,\ell}_{ij}(y)\mathfrak M_\ell(I,J\given i,j,\gamma)
C^R_\ell(y,d\gamma)\mathbb M(dy).
\label{eq:full-graph-unconditional-edge-pushforward}
\end{equation}
All three coarse event laws sum to one. Conditional endpoint independence permits the
product in \eqref{eq:full-graph-edge-pushforward} but not the factoring of the
expectation of a random event mass and a random assignment, so
\eqref{eq:full-graph-unconditional-edge-pushforward} applies to all three input laws
with no independence assumed between event and membership marks. The random
$\eta_{\ell+1}(y,\gamma)$ is part of $\mathfrak N_{\ell+1}$ with rows recovered by
\eqref{eq:full-graph-edge-event-disintegration}, and its unconditional mean does not
replace it. Loop copies arise at every depth after the first link lift
\eqref{eq:full-graph-coarse-links}; they are induced by blocking and are not a new
support declaration, as in \Cref{prop:rg-meta-attention-pushforward}. For a hard
partition the loop and product
terms agree, and under the posterior bridge the pushforward is
\eqref{eq:rg-meta-attention}. Literal unit-weight membership in several parents has
total mass greater than one; it is finite-measure replication, not a probability
pushforward. \status{ESTABLISHED}

\noindent\emph{Proof.}
Summing the first term of \eqref{eq:full-graph-edge-pushforward} over $I,J$ gives
$\sum_{i\ne j}\eta_{ij}$ and the loop term gives $\sum_i\eta_{ii}$, which add to one;
the correlated and unconditional cases use normalization of $\mathfrak M_\ell$,
$\eta_\ell$, $C^R_\ell$, and $\mathbb M$. The expectation of a product of dependent
random quantities is not the product of expectations. A vertex with unit weight in $k$
parents contributes its mass $k$ times. $\square$

\definitionheading{Chapman--Kolmogorov composition}{def:full-graph-composition}
Nested assignments compose by
\begin{equation}
\mathfrak m_{\ell+2\leftarrow\ell}(A\given i,R_\ell,R_{\ell+1})
=\sum_I\mathfrak m_{\ell+1}(A\given I,R_{\ell+1})\mathfrak m_\ell(I\given i,R_\ell),
\label{eq:full-graph-membership-composition}
\end{equation}
\begin{equation}
\mathfrak M_{\ell+2\leftarrow\ell}(A,A'\given i,j,R_\ell,R_{\ell+1})
=\sum_{I,J}\mathfrak M_{\ell+1}(A,A'\given I,J,R_{\ell+1})\mathfrak M_\ell(I,J\given i,j,R_\ell),
\label{eq:full-graph-joint-membership-composition}
\end{equation}
and the augmented structural channels by
\begin{equation}
C^{\mathrm{struct}}_{\ell+2\leftarrow\ell}
=C^{\mathrm{struct}}_{\ell+2\leftarrow\ell+1}\circ C^{\mathrm{struct}}_{\ell+1\leftarrow\ell},
\label{eq:full-graph-structural-channel-composition}
\end{equation}
where composition integrates every intermediate assignment, event law, and mark, and
the root, holonomy, and boundary marks stay in the target of each composite. An
independently declared multi-depth channel is admitted only when it coincides with the
composite, and the composite is the exact multi-depth step only when the augmented
state carries every mark that the next structural kernel reads.

\section{Coarse links and what the parent retains}
\label{sec:full-graph-vfe-nonflat}

The links are not assumed flat: $\mathfrak h_\ell$ keeps the root-framed datum of
\Cref{prop:rg-rooted-datum} in every connected block with its joint law, and on a
non-flat graph a reduced mean of the dressed marks must carry a declared residual,
whose penalization to zero is a modeling preference, not a consequence of gauge
covariance.

\definitionheading{Coarse links under a hard partition}{def:full-graph-coarse-links}
Let the realized mark $R_\ell=\gamma$ be a hard partition whose blocks are replaced by
their component-indexed meta-agents \eqref{eq:rg-component-meta-index}, so that each
coarse vertex $I$ has one root $r_I$ and one rooted spanning tree. This excludes
overlapping membership kernels, for which a copy would need a root transport into
every block that contains one of its endpoints with positive membership; the event
pushforwards of \Cref{prop:full-graph-event-pushforward} hold for general membership
kernels. The depth-$(\ell+1)$ multigraph $\mathfrak I_{\ell+1}=(V_{\ell+1},E_{\ell+1},\bar{\ })$
has the coarse vertices, the same copies $E_{\ell+1}:=E_\ell$ with the same reversal,
and endpoints $e:J\to I$ whenever $e:j\to i$ with $j\in J$ and $i\in I$; its event law
is $\eta^{x,\ell+1}_e:=\eta^{x,\ell}_e$, and its links are the dressed marks of all
copies,
\begin{equation}
\Theta^{x,\ell+1}_e
:=\widehat\Theta^{x,\ell}_e
=\Theta^{x,\ell}_{I\leftarrow i}\Theta^{x,\ell}_e\bigl(\Theta^{x,\ell}_{J\leftarrow j}\bigr)^{-1},
\qquad e:j\to i,\ j\in J,\ i\in I,
\label{eq:full-graph-coarse-links}
\end{equation}
internal as well as boundary, so that internal copies become loops at their block. The
parent laws $q^x_{I,\ell+1}$ are expressed in the frame of the root $r_I$. The coarse
multigraph is induced structure; no edge is created.

\propositionheading{Properties of the coarse links}{prop:full-graph-coarse-links}
The coarse links are well typed and keep the whole link content of depth $\ell$. An
internal tree copy has the identity link $1_G$. An internal non-tree copy carries the
based holonomy of its fundamental loop, and for every closed walk $p$ in $I$ based at
$r_I$ that traverses the copies $e_0,\ldots,e_{n-1}$ in this order,
$\mathrm{Hol}^{x,\ell}_I(p)=\Theta^{x,\ell+1}_{e_{n-1}}\cdots\Theta^{x,\ell+1}_{e_0}$.
Inverse reversal persists, and the dressed law is the vertex law
\eqref{eq:geo-regime-two-gauge-law} with vertex coordinates $g_{r_I}$, so the
divergences, rows, and row free energies of
\Cref{def:full-graph-transported-divergences} are defined and gauge invariant at
depth $\ell+1$. On a loop the divergence
$\KL(q^x_{I,\ell+1}\Vert(\Theta^{x,\ell+1}_e)_\#q^x_{I,\ell+1})$ vanishes on tree copies
and, on a non-tree copy, exactly when the parent law is fixed by that holonomy, as in
\eqref{eq:full-graph-holonomy-stabilization}. The pair marginal
$\sum_{e:J\to I}\eta^{x,\ell+1}_e$ is \eqref{eq:full-graph-edge-pushforward} for the
hard partition, and no averaged group element replaces the parallel copies. The links
depend on the trees, not only on the roots: changing the root of a fixed tree is a
vertex gauge transformation at depth $\ell+1$, while replacing the tree changes
$\Theta^{x,\ell}_{I\leftarrow i}$ to $\mathrm{Hol}^{x,\ell}_I(p_i)\Theta^{x,\ell}_{I\leftarrow i}$
for a based loop $p_i$ that depends on $i$, and a copy that leaves the tree changes its
loop link from $1_G$ to the holonomy $\mathrm{Hol}^{x,\ell}_I(p)$ of its new fundamental
loop $p$ and its loop divergence from $0$ to
$\KL(q^x_{I,\ell+1}\Vert\mathrm{Hol}^{x,\ell}_I(p)_\#q^x_{I,\ell+1})$. Whenever such a
holonomy is not the identity, the trees are part of the coarse channel rather than a
gauge choice. When $q^x_{I,\ell+1}$ and $q^x_{J,\ell+1}$ lie in the stabilized sector,
the divergences $D^{x,\ell+1}_e$ do not depend on the trees. This link-level tree
dependence complements the feature-level one of
\eqref{eq:rg-linear-cross-scale-tree-change}. \status{ESTABLISHED}

\noindent\emph{Proof.}
The dressed link of an internal copy is $1_G$ on a tree copy, since there
$\Theta_e=\Theta_{I\leftarrow i}^{-1}\Theta_{I\leftarrow j}$, and is the
fundamental-cycle holonomy \eqref{eq:full-graph-tree-transport-law} otherwise.
Inserting $(\Theta_{I\leftarrow k})^{-1}\Theta_{I\leftarrow k}$ at each intermediate
vertex $k$ of a closed walk turns the ordered product of the fine links into that of
the coarse links. Reversal and the vertex law follow from
\eqref{eq:rg-tree-transport-law}, and relative entropy is invariant under a common
bijective pushforward, which gives gauge invariance; a loop divergence vanishes
exactly when $q_I=\mathrm{Hol}_I(p)_\#q_I$. A root change multiplies every transport
of the block by one element, and a tree change multiplies each tree transport on the
left by a based holonomy (\Cref{prop:rg-rooted-datum}), so the link of $e:j\to i$ becomes
$\mathrm{Hol}_I(p_i)\Theta^{x,\ell+1}_e\mathrm{Hol}_J(p_j)^{-1}$, and when
$\mathrm{Hol}_I(p_i)$ fixes $q_I$ and $\mathrm{Hol}_J(p_j)$ fixes $q_J$, invariance under
the pushforward by $\mathrm{Hol}_I(p_i)^{-1}$ gives
$\KL(q_I\Vert(\mathrm{Hol}_I(p_i)\Theta\mathrm{Hol}_J(p_j)^{-1})_\#q_J)=\KL(q_I\Vert\Theta_\#q_J)$. $\square$

\section{How a parent influences its children}
\label{sec:full-graph-vfe-downward}

\propositionheading{Downward coupling in a declared hierarchy}{prop:full-graph-downward-coupling}
Under \Cref{hyp:full-graph-depth-markov} for a declared top-down hierarchy, a
factorization such as
\begin{equation}
K^\ell_\downarrow(dZ_\ell\given Y_{\ell+1},R_\ell,\mathfrak N_\ell,X_*)
=\prod_{I\in V_{\ell+1}}K^\ell_{\downarrow,I}\bigl(dZ_{I,\ell}\given Z_{I,\ell+1},\mathfrak N_{\ell+1},R_\ell,\mathfrak N_\ell,X_*\bigr)
\label{eq:full-graph-downward-kernel}
\end{equation}
makes the parent belief and model presentations arguments of the children's
conditional generative prior, and the corresponding term of
\eqref{eq:full-graph-vfe-factorized} is
\begin{equation}
\Fenergy^\ell_{\mathrm{cross}}
=\E_{\mathbb Q_\lambda}
\KL\bigl(\mathbb Q^\ell_Z(\cdot\given Y_{\ell+1},R_\ell,\mathfrak N_\ell,o,X_*)\Vert
K^\ell_\downarrow(\cdot\given Y_{\ell+1},R_\ell,\mathfrak N_\ell,X_*)\bigr).
\label{eq:full-graph-cross-scale-kl}
\end{equation}
The generative denominator in this cross-scale divergence reads the parent coordinates through
$K^\ell_\downarrow$. The full tower objective can also couple parent coordinates to child structural
marks through $\mathbb P^\ell_R$ and $\mathbb P^\ell_N$ in \Cref{hyp:full-graph-depth-markov}.
A downward kernel that does not depend on the parent coordinates
contributes no parent--child cross-derivative through this generative factor, and
normalization of the downward kernel alone implies no coupling. When
$K^\ell_\downarrow$ depends on the parent coordinates and the recognition
parameterization transmits that dependence, the mixed derivative in a parent and a
child parameter can be nonzero, so that optimization is coupled in both directions
although the generative arrow is top-down; a universal claim of such coupling for a
family requires a proved nonzero cross-derivative in that family. For a binary parent
$b$ and child $y$ with product recognition factors $q_b(1)=u$ and $q_y(1)=v$, the
mixed derivative of $\E_{q_bq_y}[-\log K_\downarrow(y\given b)]$ in $(u,v)$ is
$-\log[K_\downarrow(1\given1)K_\downarrow(0\given0)/(K_\downarrow(0\given1)K_\downarrow(1\given0))]$,
which vanishes exactly when the odds ratio of the kernel is one.
\status{ESTABLISHED}

\noindent\emph{Proof.}
The generative denominator in \eqref{eq:full-graph-cross-scale-kl} is $K^\ell_\downarrow$.
If that kernel does not read the parent coordinates, its negative log density has no parent
argument; the expectation and recognition parameterization may still depend on the parent. For the binary case the expectation is bilinear in $(u,v)$, and its $uv$
coefficient is the displayed alternating sum of $-\log K_\downarrow(y\given b)$ over
the four configurations. $\square$

This is the static sense in which a meta-agent influences its sub-agents without a
generative factor reading a posterior; in an adaptive tower the same kernel is the
reversal kernel of \Cref{lem:full-graph-adaptive-reversal}.

\section{Closing the tower}
\label{sec:full-graph-closing-tower}

A tower in which each depth's prior is supplied by the depth above must terminate, and
within one finite acyclic joint its top cannot be constituted by the agents it
organizes (\Cref{prop:obs-declared-root-unavoidable}). A same-time reciprocal
structure among latents is admissible only as a globally normalized joint
(\Cref{sec:gen-undirected}, \Cref{prop:gen-gibbs-counterexample}): a flat unanchored
energy is improper, an anchored or non-flat energy with its global normalizer is a
valid undirected model, and an energy appended without its normalizer is not an
evidence bound. A same-time deterministic parent $Z_{\ell+1}=\chi_\ell(Z_\ell)$ and an
independently sampled parent that generates $Z_\ell$ are therefore not two free
descriptions of one joint: one uses the hierarchy above, a deterministic pushforward,
or a normalized undirected compatibility factor, and reciprocal same-time influence
requires the last with its partition function. Supplement G works out the anchored
Gaussian star and the convergence of its coordinate ascent.

\subsection{A declared root}
\label{sec:obs-survives}

A truncated tower with a declared top prior $\mathbb P_L$ satisfies the hypotheses of
\Cref{ch:elbo} without amendment; its top is given, not constituted.

\propositionheading{A declared root is unavoidable}{prop:obs-declared-root-unavoidable}
Let a finite directed acyclic generative model be typed as in \Cref{ch:generative}, so
that every observed record is a child of at least one latent vertex. Then at least one
latent has no parents, and its factor in the joint depends only on the parameters and
the structural data. No finite acyclic construction of this type has every prior
constituted by other latents. \status{ESTABLISHED}

\noindent\emph{Proof.}
A finite directed acyclic graph admits a topological order, whose first element has
no incoming edge. That element is not a record, because every record has a latent
parent, so it is a latent, and its factor is conditioned on no other vertex. The typing
hypothesis is needed: in the graph $o\to y$ with observed root $o$, the only latent has
a parent. $\square$

Anchored, globally normalized cyclic models remain available, but their anchors and
partition functions are declared; parameter learning fits only their parameters.

\subsection{Three loci of participation}
\label{sec:obs-participatory}

What a cycle in the model was reaching for is partly available without one, at three
loci, each a sequence of evaluations of fixed normalized joints.

\remarkheading{The inference loop is block coordinate ascent}{rem:full-graph-inference-loop}
In a mean-field recognition family, the coordinate update of a factor reads the current
factors of every variable that shares a factor of the joint with it
\citep[Sec.~10.1.1]{Bishop2006}. In every connected hierarchical model in which the
children's conditionals depend on their parents, parent and child updates therefore
read each other; at a stationary point of the schedule the factor that acts as a child's
prior in its own update is $\exp\E[\log K_\downarrow(z_{\mathrm{child}}\given
\cdot)]$ under the parent's current factor, which is itself formed from the factors of
all its children. The effective prior is thus formed from the population's current
beliefs through a latent inside the joint. This is a schedule of updates against one
fixed, acyclic, normalized joint, not participation in time, and no kernel reads a
posterior, so \Cref{req:gen-typing-prohibition} holds. \status{STANDARD}

\propositionheading{The aggregate recognition law is the optimal top prior}{prop:full-graph-aggregate-prior}
Let episodes $n=1,\ldots,N_{\mathrm e}$ be conditionally independent given $\theta$,
each with its own top latent in a common space carrying the prior density $p_\theta$
relative to a fixed reference measure, let $\mathbb Q_n$ be the recognition law of the
top latent of episode $n$, and let every other factor be free of $\theta$. The
$\theta$-dependent part of the summed bound is
\begin{equation}
\sum_{n=1}^{N_{\mathrm e}}\E_{\mathbb Q_n}[\log p_\theta]
=N_{\mathrm e}\E_{\bar{\mathbb Q}}[\log p_\theta],
\qquad
\bar{\mathbb Q}=\frac1{N_{\mathrm e}}\sum_{n=1}^{N_{\mathrm e}}\mathbb Q_n .
\label{eq:obs-aggregate-prior}
\end{equation}
If $\bar{\mathbb Q}$ has a density $\bar q$ with finite entropy relative to the same
reference, then $\E_{\bar{\mathbb Q}}[\log p]=\E_{\bar{\mathbb Q}}[\log\bar
q]-\KL(\bar{\mathbb Q}\Vert p)$ for every prior density $p$. Over an unrestricted
family of prior densities the bound is maximized exactly at the aggregate recognition
law $p=\bar q$, and within a restricted family its maximizers are the minimizers of
$\KL(\bar{\mathbb Q}\Vert p)$ over the family, when they exist. Each evaluation uses one
fixed normalized joint, and updating $\theta$ at fixed recognition is a coordinate
step of the two-argument objective \citep{Neal1998}, so
\Cref{req:gen-typing-prohibition} holds. \status{ESTABLISHED}

\noindent\emph{Proof.}
The first identity is linearity of expectation in the recognition law. The second is
$\log p=\log\bar q-\log(\bar q/p)$ integrated against $\bar{\mathbb Q}$, and Gibbs'
inequality gives the maximizers. $\square$

This is the variational counterpart of empirical Bayes \citep{Robbins1956}, whose
optimum is the aggregate-posterior prior of \citet{HoffmanJohnson2016} and
\citet{TomczakWelling2018}; only the parameter of the declared prior family is fitted,
by a slow-block update in the typing of \Cref{sec:obs-fast-slow}.

\definitionheading{Episode-indexed joint}{def:full-graph-episode-joint}
Episodes are indexed by $n=0,1,\ldots$ with records $o_n$. In the parameter route each
episode has its own normalized joint $\mathbb P_{\theta_n,V}$, and the parameter is
updated by
\begin{equation}
\theta_{n+1}=\mathcal U_n\bigl(\theta_n,\mathbb Q^{(n)}_{V,o_n,X}\bigr),
\label{eq:obs-episode-update}
\end{equation}
where $\mathbb Q^{(n)}_{V,o_n,X}$ is the recognition coupling selected at episode $n$
and $\mathcal U_n$ is a declared map; within each episode the joint is fixed, and its
generative kernels read $\theta_n$ and never a recognition law of the same episode. In
the latent route one joint spans episodes $0,\ldots,n_1$: episode $n$ has a tower state
$\mathbf Y^{(n)}$ with parent coordinates $Z^{(n)}_{\mathrm{par}}=\chi^{\mathrm{ep}}(\mathbf Y^{(n)})$
for a measurable read-out $\chi^{\mathrm{ep}}$, and
\begin{equation}
\begin{aligned}
\mathbb P^{\mathrm{ep}}\bigl(d\mathbf Y^{(0:n_1)},do_{0:n_1}\given X\bigr)
={}&\mathbb P^{(0)}(d\mathbf Y^{(0)}\given X)
\prod_{n=1}^{n_1}\mathbb P^{(n\leftarrow n-1)}\bigl(d\mathbf Y^{(n)}\given Z^{(n-1)}_{\mathrm{par}},X\bigr)\\
&\times\prod_{n=0}^{n_1}\mathbb P^{\mathrm{obs}}\bigl(do_n\given\mathbf Y^{(n)},X\bigr),
\end{aligned}
\label{eq:full-graph-episode-joint}
\end{equation}
with a declared root in $\mathbb P^{(0)}$; the parent latent of episode $n-1$
parameterizes the conditional law of the tower of episode $n$. The observation
kernel is any Markov kernel and need not have a density: the tower state may carry
a latent copy $\widetilde o_n$ of the record, observed through
$\mathbb P^{\mathrm{obs}}(do_n\given\mathbf Y^{(n)},X)=\delta_{\widetilde o_n}(do_n)$,
in which case the free energy of the joint is its extended form
\eqref{eq:full-graph-vfe-global} and the decomposition
\eqref{eq:full-graph-episode-decomposition} applies as displayed only for discrete
records.

\propositionheading{Delayed feedback in the episode-indexed joint}{prop:full-graph-episode-joint}
In the latent route, \eqref{eq:full-graph-episode-joint} is a normalized joint whose
kernels read structural data and earlier latents only, so
\Cref{req:gen-typing-prohibition} holds; its graph is acyclic with every
cross-episode edge pointing forward in $n$, so
\Cref{prop:obs-declared-root-unavoidable} still requires a declared root in episode
$0$. The prior of the tower at episode $n+1$ given the records so far is the filtered
parent law pushed through the transition,
\begin{equation}
\mathbb P^{\mathrm{ep}}\bigl(\mathbf Y^{(n+1)}\in\cdot\given o_{0:n}\bigr)
=\int\mathbb P^{(n+1\leftarrow n)}(\cdot\given z,X)
\mathbb P^{\mathrm{ep}}\bigl(Z^{(n)}_{\mathrm{par}}\in dz\given o_{0:n}\bigr),
\label{eq:full-graph-episode-predictive}
\end{equation}
so the children's prior at episode $n+1$ is set by the parent inferred from the
population at episode $n$. For every recognition law $\mathbb Q$ on $\mathbf
Y^{(0:n_1)}$, on the finite-density domain with observation densities
$p^{\mathrm{obs}}$,
\begin{equation}
\begin{aligned}
\Fenergy^{\mathrm{ep}}[\mathbb Q]
={}&\sum_{n=0}^{n_1}\E_{\mathbb Q}\bigl[-\log p^{\mathrm{obs}}(o_n\given\mathbf Y^{(n)},X)\bigr]
+\KL\bigl(\mathbb Q^{(0)}\Vert\mathbb P^{(0)}\bigr)\\
&+\sum_{n=1}^{n_1}\E_{\mathbb Q}\KL\bigl(\mathbb Q(\cdot\given\mathbf Y^{(0:n-1)})\Vert
\mathbb P^{(n\leftarrow n-1)}(\cdot\given Z^{(n-1)}_{\mathrm{par}},X)\bigr),
\end{aligned}
\label{eq:full-graph-episode-decomposition}
\end{equation}
where $\mathbb Q^{(0)}$ is the law of $\mathbf Y^{(0)}$ under $\mathbb Q$ and
$\mathbb Q(\cdot\given\mathbf Y^{(0:n-1)})$ the conditional law of $\mathbf Y^{(n)}$.
In the parameter route, if $\mathcal U_n$ maximizes the bound of episode $n$ over the
parameter at fixed $\mathbb Q^{(n)}_{V,o_n,X}$, the bound of episode $n$ at
$\theta_{n+1}$ is at least its value at $\theta_n$. \status{ESTABLISHED}

\noindent\emph{Proof.}
Normalization is composition of Markov kernels, as in
\Cref{prop:gen-exact-normalization}, and acyclicity follows because every kernel reads
only earlier episodes. Given $Z^{(n)}_{\mathrm{par}}$, the tower of episode $n+1$ is
independent of $(\mathbf Y^{(0:n)},o_{0:n})$, because its kernel reads only
$Z^{(n)}_{\mathrm{par}}$ and the records up to $n$ are generated before it; conditioning
on $o_{0:n}$ and integrating out $Z^{(n)}_{\mathrm{par}}$ gives
\eqref{eq:full-graph-episode-predictive}. The decomposition is the chain rule
\Cref{lem:prob-kl-chain-rule} applied along $n$ to the Radon--Nikodym ratio of
$\mathbb Q$ against the joint. The last statement is the definition of a maximizer.
$\square$

\propositionheading{One joint for a declared parameter route}{prop:full-graph-episode-lifted-joint}
Let the episode joints of the parameter route form a Markov kernel
$\theta\mapsto\mathbb P_{\theta,V}(dy,do\given X)$ on the parameter space, with
evidence densities $p_{\theta,V}(o\given X)$ relative to one reference $\nu$ on
$\mathsf O$ and posterior slices $\boldsymbol\Pi_{\theta,V,o,X}$ jointly measurable in
$(\theta,o)$; let the coupling $\mathbb Q^{\theta}_{V,o,X}$ selected by
\Cref{hyp:prob-measurable-coupling-selector} for the joint at $\theta$, so that
$\mathbb Q^{(n)}_{V,o_n,X}=\mathbb Q^{\theta_n}_{V,o_n,X}$, be jointly measurable in
$(\theta,o)$; and let each update $\mathcal U_n$ of
\eqref{eq:obs-episode-update} be measurable, so that
$\mathcal U^{\mathrm{sel}}_n(\theta,o):=\mathcal U_n(\theta,\mathbb Q^{\theta}_{V,o,X})$
is measurable and the route reads $\theta_{n+1}=\mathcal U^{\mathrm{sel}}_n(\theta_n,o_n)$
from a declared $\theta_0$. Write $\theta_n(o_{0:n-1})$ for the parameter the route
reaches after the records $o_{0:n-1}$, and let each realized record be admitted at
that parameter. The lifted joint on the tower states
$\mathbf Y^{(n)}=(\theta_n,y_n,\widetilde o_n)$, $n=0,\ldots,n_1$, and the records,
\begin{equation}
\mathbb P^{\mathrm{lift}}\bigl(d\mathbf Y^{(0:n_1)},do_{0:n_1}\given X\bigr)
=\delta_{\theta_0}(d\theta_0)
\prod_{n=0}^{n_1}\mathbb P_{\theta_n,V}(dy_n,d\widetilde o_n\given X)\delta_{\widetilde o_n}(do_n)
\prod_{n=1}^{n_1}\delta_{\mathcal U^{\mathrm{sel}}_{n-1}(\theta_{n-1},\widetilde o_{n-1})}(d\theta_n),
\label{eq:full-graph-lifted-joint}
\end{equation}
has four properties. (i) It is a latent route of \Cref{def:full-graph-episode-joint}
with parent coordinates $Z^{(n)}_{\mathrm{par}}=(\theta_n,\widetilde o_n)$, transition
$\mathbb P^{(n\leftarrow n-1)}(d\mathbf Y^{(n)}\given z,X)
=\delta_{\mathcal U^{\mathrm{sel}}_{n-1}(z)}(d\theta_n)\mathbb P_{\theta_n,V}(dy_n,d\widetilde o_n\given X)$,
observation kernel $\delta_{\widetilde o_n}(do_n)$, and root
$\delta_{\theta_0}\otimes\mathbb P_{\theta_0,V}$, so \Cref{prop:full-graph-episode-joint}
applies to it. (ii) Given $o_{0:n-1}$, the conditional law of $(y_n,\widetilde o_n)$ is
the episode joint $\mathbb P_{\theta_n(o_{0:n-1}),V}$ of the parameter route, and the
evidence density of the lifted joint relative to $\nu^{\otimes(n_1+1)}$ is the product
of the episode evidences,
\begin{equation}
p^{\mathrm{lift}}(o_{0:n_1}\given X)=\prod_{n=0}^{n_1}p_{\theta_n(o_{0:n-1}),V}(o_n\given X).
\label{eq:full-graph-lifted-evidence}
\end{equation}
(iii) Its posterior slice fixes $\widetilde o_n=o_n$ and $\theta_n=\theta_n(o_{0:n-1})$
and is the product $\bigotimes_n\boldsymbol\Pi_{\theta_n(o_{0:n-1}),V,o_n,X}$ on
$y_{0:n_1}$, so the posterior law of $y_n$ reads $o_{0:n}$ only. (iv) For every
recognition law $\mathbb Q$ on the tower states, the free energy
\eqref{eq:full-graph-vfe-global} of the lifted joint is
\begin{equation}
\Fenergy^{\mathrm{lift}}[\mathbb Q]
=\sum_{n=0}^{n_1}\Bigl[-\log p_{\theta_n,V}(o_n\given X)
+\KL\bigl(\mathbb Q^{y}_n\Vert\boldsymbol\Pi_{\theta_n,V,o_n,X}\bigr)\Bigr]
+\operatorname{TC}(\mathbb Q^{y}),
\qquad\theta_n=\theta_n(o_{0:n-1}),
\label{eq:full-graph-lifted-free-energy}
\end{equation}
when $\mathbb Q$ fixes $\widetilde o_n=o_n$ and $\theta_n=\theta_n(o_{0:n-1})$ almost
surely, and $+\infty$ otherwise, where $\mathbb Q^{y}$ is the law of $y_{0:n_1}$ under
$\mathbb Q$, $\mathbb Q^{y}_n$ its marginal on $y_n$, each bracket is the extended free
energy \eqref{eq:elbo-extended-free-energy} of episode $n$ at the parameter the route
reaches, and $\operatorname{TC}$ is the total correlation
\eqref{eq:obs-total-correlation} across episodes. At product laws, and in particular at
the selected couplings $\bigotimes_n\mathbb Q^{\theta_n}_{V,o_n,X}$, the free energy
is the sum of the episode bounds of the parameter route. \status{ESTABLISHED}

\noindent\emph{Proof.}
Each factor of \eqref{eq:full-graph-lifted-joint} is a Markov kernel in the
coordinates to its left: $\theta\mapsto\mathbb P_{\theta,V}$ by hypothesis, and the
Dirac factors as kernels of measurable maps, $\mathcal U^{\mathrm{sel}}_{n-1}$ being
the composition of $\mathcal U_{n-1}$ with the jointly measurable selection. Finite
composition normalizes the product (\Cref{prop:gen-exact-normalization}), and reading
the factors in the order of \eqref{eq:full-graph-episode-joint} gives (i); every kernel
reads structural data and latent coordinates only, because the selection is evaluated
at the latent copy $\widetilde o_{n-1}$ and not at a recognition law of the joint being
scored. For (ii), integrate \eqref{eq:full-graph-lifted-joint} against a bounded
measurable $f(o_{0:n-1},y_n,\widetilde o_n)$: the factors of episodes $n+1,\ldots,n_1$
have unit mass, the Dirac factors identify $o_k=\widetilde o_k$ for $k<n$ and, by
induction along the recursion, $\theta_n=\theta_n(o_{0:n-1})$, and Tonelli's theorem
leaves the integral of
$\int f(o_{0:n-1},y,\widetilde o)\mathbb P_{\theta_n(o_{0:n-1}),V}(dy,d\widetilde o\given X)$
against the law of $o_{0:n-1}$, which is the stated conditional law. Its record
marginal is $p_{\theta_n(o_{0:n-1}),V}(o_n\given X)\nu(do_n)$, and multiplying the
conditional densities along $n$ gives \eqref{eq:full-graph-lifted-evidence}; the
product is measurable because the evidence density is jointly measurable and
$o_{0:n-1}\mapsto\theta_n(o_{0:n-1})$ is a composition of measurable maps, and it is
finite and positive because each factor is. For (iii), substitute the disintegration
$\mathbb P_{\theta,V}(dy,d\widetilde o\given X)
=\boldsymbol\Pi_{\theta,V,\widetilde o,X}(dy)p_{\theta,V}(\widetilde o\given X)\nu(d\widetilde o)$
into each factor and integrate out the Dirac coordinates: for bounded measurable
$f$ of the tower states and $g$ of the records,
\begin{equation}
\begin{aligned}
\E_{\mathbb P^{\mathrm{lift}}}[fg]
={}&\int g(o)\Bigl[\int f\bigl((\theta_n(o_{0:n-1}),y_n,o_n)_{n\le n_1}\bigr)
\bigotimes_{n=0}^{n_1}\boldsymbol\Pi_{\theta_n(o_{0:n-1}),V,o_n,X}(dy_n)\Bigr]\\
&\times p^{\mathrm{lift}}(o\given X)\nu^{\otimes(n_1+1)}(do),
\end{aligned}
\label{eq:full-graph-lifted-posterior}
\end{equation}
and the bracket, measurable in $o$ by the joint measurability of the slices, is a
version of the conditional expectation of $f$ given the records. For (iv), the
posterior is carried by the set on which $\widetilde o_n=o_n$ and
$\theta_n=\theta_n(o_{0:n-1})$ for every $n$; a $\mathbb Q$ that charges its
complement is not absolutely continuous with respect to the posterior, and its
relative entropy is $+\infty$. Otherwise disintegrate both laws over $y_{0:n_1}$ with
\Cref{lem:prob-kl-chain-rule}(a): the conditional laws of the remaining coordinates
given $y_{0:n_1}$ are the same Dirac mass under both, so
$\KL(\mathbb Q\Vert\boldsymbol\Pi^{\mathrm{lift}})
=\KL(\mathbb Q^{y}\Vert\bigotimes_n\boldsymbol\Pi_{\theta_n,V,o_n,X})$, and the
product-reference identity \eqref{eq:elbo-total-correlation-chain} splits the right
side into $\operatorname{TC}(\mathbb Q^{y})+\sum_n\KL(\mathbb Q^{y}_n\Vert\boldsymbol\Pi_{\theta_n,V,o_n,X})$.
Adding $-\log p^{\mathrm{lift}}(o\given X)$ from \eqref{eq:full-graph-lifted-evidence}
gives \eqref{eq:full-graph-lifted-free-energy}, and the total correlation of a product
law is zero. $\square$

The route's bookkeeping with the parameter read from the law being scored is the
composite that \Cref{req:gen-typing-prohibition} excludes; the lifted joint evaluates
the selection at the latent copy of the record instead, and its evidence is the
sequential empirical-Bayes evidence of the route.

\openproblemheading{Accounting across episodes}{open:full-graph-episode-accounting}
In the parameter route the bounds of different episodes refer to different joints and
records. Within an episode a maximizing update is monotone
(\Cref{prop:full-graph-episode-joint}), and for an update that reads a coupling
selected by a declared rule the lifted joint of
\Cref{prop:full-graph-episode-lifted-joint} is one joint whose evidence is the product
of the episode evidences and whose free energy is the sum of the episode bounds plus
the cross-episode total correlation; an update that reads the recognition law being
scored is excluded by \Cref{req:gen-typing-prohibition}. The lifted joint informs the
parent of episode $n+1$ by the record of episode $n$ through its latent copy. What is
unproved is whether the label coordinates alone can carry the feedback. The problem is
posed for label-sufficient latent routes of \Cref{def:full-graph-episode-joint}: the
episode model, with a finite label coordinate $y^{\mathrm{lab}}$ whose prior is the
parameter, its private latents, and its record kernel, is that of the parameter route
with the posterior slice as its recognition coupling and the aggregate update of
constant step,
$\theta_{n+1}=(1-\epsilon_{\mathrm e})\theta_n+\epsilon_{\mathrm e}\boldsymbol\Pi^{\mathrm{lab}}_{\theta_n,V,o_n,X}$
with $\boldsymbol\Pi^{\mathrm{lab}}$ the label marginal of the posterior slice; the
tower state of episode $n$ consists of the episode model's latents and an auxiliary
preference coordinate; each record is emitted from its own episode's latents by the
record kernel; and the read-out $\chi^{\mathrm{ep}}$ reads the label and the
preference coordinate only. The records either have one value whose density
$p_V(o\given y^{\mathrm{lab}},X)$ given the label does not depend on the label, while
the expected squared total-variation displacement of the label posterior from the
prior, under the evidence at that prior, is bounded below uniformly over priors whose
smallest weight is at least half that of $\theta_0$; or their likelihood vectors
$(p_V(o\given y^{\mathrm{lab}},X))_{y^{\mathrm{lab}}}$ span a cone on which no nonzero
polynomial vanishes. For such routes, with the per-episode discrepancy the total
variation between the route's filtered label law
$\mathbb P^{\mathrm{ep}}(y^{\mathrm{lab}}_{n+1}\in\cdot\given o_{0:n})$ and
$(1-\epsilon_{\mathrm e})\mathbb P^{\mathrm{ep}}(y^{\mathrm{lab}}_{n}\in\cdot\given o_{0:n-1})
+\epsilon_{\mathrm e}\mathbb P^{\mathrm{ep}}(y^{\mathrm{lab}}_{n}\in\cdot\given o_{0:n})$,
the aggregate update applied to the route's own filtered law and label posterior, the
problem is settled by a route, defined for every step, whose discrepancy is
$o(\epsilon_{\mathrm e})$ uniformly over record histories and over episodes $n$ up to
any fixed multiple of $1/\epsilon_{\mathrm e}$, or by a lower bound of order
$\epsilon_{\mathrm e}$ on the supremum of the discrepancy over such histories and
episodes for every such route. Whether a declared update intertwines the fine and
parent evolutions is part of \Cref{open:rg-pointwise-parent-dynamics}; $\theta_n$
belongs to the slow block of \Cref{def:obs-fast-slow-blocks}.

\subsection{The remaining hierarchy question}
\label{sec:full-graph-vfe-target}

The prepared parent of \Cref{sec:rg-memory-parent-mechanism}, the memory cost of
partial refresh in \Cref{sec:rg-memory-projection}, and the successive formation of
\Cref{thm:finite-depth-adaptive-formation} are constructions at stated hypotheses, not
consequences of arbitrary descent.

\openproblemheading{Persistent multiscale hierarchy}{open:full-graph-hierarchy}
Starting from $c_*$ with a finite multigraph datum $\mathfrak N_0$, construct a
normalized partition law and parent--child kernels such that descent of
\eqref{eq:full-graph-vfe-global} produces persistent nested blocks, the pushforward
rules \eqref{eq:full-graph-edge-pushforward} to
\eqref{eq:full-graph-structural-channel-composition} and the link lift
\eqref{eq:full-graph-coarse-links} compose across depths, the parent family has a
controlled closure residual in the sense of \Cref{prop:rg-interaction-residual}, and
non-flat holonomy is kept or stabilized without imposing flatness. What holds: every
candidate hierarchy has the exact accounting of
\Cref{thm:full-graph-vfe-decomposition}; at one depth the fast--slow typing of
\Cref{def:obs-fast-slow-blocks} and its conditional slow-manifold limit,
\Cref{open:obs-slow-manifold}, are available; \Cref{thm:finite-depth-adaptive-formation}
forms nested blocks under a supplied energy hierarchy whose strengths are tied to the
temperature; and \Cref{prop:rg-formation-endogenous-cycle} forms one level. What is
unproved is that descent produces a second level that depends on the first. The
problem is posed in the class of models at $c_*$ whose structural marks are drawn from
a finite partition catalog with a partition factor inside the generative and the
recognition law (\Cref{def:rg-geometry-ensemble,def:full-graph-partition-factor}),
whose recognition coordinates descend \eqref{eq:full-graph-vfe-global} by a reversible
diffusion at temperature $\mathrm T$ with rate-one redraws of the discrete marks from
their conditional laws, as in \Cref{def:rg-formation-cycle-experiment}, and whose
energies, precisions, prior ratios, hyperparameters, and mobilities do not depend on
$\mathrm T$, as $\mathrm T\to0$. It is settled by three results together, or by a
proof that the first of them cannot hold in this class. The first is a lineage of at
least two levels, formed with probability bounded below from every initial law with a
density and a finite second moment, whose second level meets two conditions. It
depends on the first level: for a declared handoff time, a declared coalescence event
whose block is a union of at least two first-level blocks and whose windows are
measured from the handoff, one declared statistic of the handoff state (the
probability of the event, of its occurrence by declared times, or of a declared
block), and a handoff class fixed before the comparison that contains every
first-level handoff and, with probability bounded below, the handoffs of histories
without first-level formation, the event has probability bounded below after
first-level formation, does not have probability tending to one uniformly over the
handoff class at the declared windows, and the statistic exceeds its value after a
counterfactual handoff in the class by a margin bounded below. It is a coherence
episode rather than restored equilibrium or coincidental coherence: a restoring
interaction, evaluated on the law of the coordinates given the realized marks or on
their marginal over the marks and not on a conditional law held fixed by hand, holds
the second-level block's coherence at a declared resolution, with the restoring rate
times the squared resolution divided by $\mathrm T$ tending to infinity, on horizons
short compared with the stationary time of entrance into the event from a set of
stationary mass bounded below, or with a probability exceeding its stationary value by
a margin bounded below. The second result is a nondegenerate intrinsic partition
selector (\Cref{open:cg-natural-selector}), and the third is a separated slow manifold
on which depth-$(\ell+1)$ variables feed back consistently into depth-$\ell$ kernels.
The negative resolution is a proof that in this class no second level produced by
descent, apart from a supplied-scale sequence whose event occurs with probability
tending to one uniformly over the handoff class at the declared windows, both meets
the restoring and horizon conditions and exceeds its counterfactual statistic by
such a margin.

\section*{Summary for physicists}

The multiscale free energy is that of one generative and one recognition law over the
whole tower, not a sum of energies chosen scale by scale. Bayes reversal of a
bottom-up structural channel conditions every depth on the whole next-depth state,
graph datum included. Partitions enter as a normalized factor, normalized branch by
branch or globally, and coarse links are the dressed links of all copies, so internal
loops keep their holonomy. A parent can couple to child coordinates through a downward kernel that
reads it and to child structural marks through the partition and graph factors. A finite acyclic tower has a declared root, and the agents
participate in their own top through coordinate ascent, a learned aggregate prior, and
delayed feedback. A sequence of episodes whose parameter update reads a coupling
selected by a declared rule is itself one joint, with the product of the episode
evidences as its evidence (\Cref{prop:full-graph-episode-lifted-joint}).
